\documentclass[aps,prfluids,reprint,onecolumn,amsmath,amssymb,nofootinbib,floatfix]{revtex4-2}

\usepackage{graphicx}
\usepackage{amsmath,amssymb}
\usepackage{booktabs}
\usepackage{xcolor}
\usepackage{hyperref}
\usepackage{enumitem}

\newcommand{\code}[1]{\texttt{\small #1}}
\setlist[itemize]{itemsep=1pt, topsep=3pt, parsep=0pt}

\begin{document}

\title{Hybrid Tensor-Train Time Evolution:\\ Status, Findings, and Outlook}

\author{S. Pisoni}
\affiliation{TN\_Turbulence project --- 2D passive scalar study}

\date{\today}

\begin{abstract}
Concise status note on the machinery built to evolve a 2D scalar field in the
hybrid tensor-train (TT) encoding of the compression study. Records the
equations being solved, what has been implemented, what has been checked, the
findings and caveats so far, and the plan going forward. Detailed numbers live
in the test suite; this note keeps only the conclusions. As of this revision the
machinery is validated, a genuine cascade has been run, and the central
question --- whether the TT rank stays bounded once fine scales are generated ---
has been answered affirmatively.
\end{abstract}

\maketitle

\section{The equations}

\subsection{Passive scalar advection--diffusion}

The target problem, and the one the compression study characterises statistically:
\begin{equation}
\partial_t \theta + \mathbf{u}\!\cdot\!\nabla\theta \;=\; \kappa\nabla^2\theta + f,
\qquad \nabla\!\cdot\!\mathbf{u} = 0,
\label{eq:ps}
\end{equation}
on a doubly periodic square. Here $\theta$ is a scalar (temperature, dye
concentration) carried passively by a prescribed incompressible velocity
$\mathbf{u}$: it does not react back on the flow. $\kappa$ is the molecular
diffusivity and $f$ a large-scale source.

\paragraph{Physics.} Advection stretches and folds $\theta$, cascading scalar
variance from the injection scale down to the Batchelor scale where diffusion
erases it. The resulting fields are strongly \emph{intermittent}: increments have
heavy tails and flatness grows without bound as the separation shrinks, which is
precisely the property the hybrid encoding was designed to preserve. The variance
budget follows from Eq.~\eqref{eq:ps},
\begin{equation}
\frac{d}{dt}\,\tfrac12\langle\theta^2\rangle
\;=\; -\,\kappa\big\langle|\nabla\theta|^2\big\rangle \;+\; \langle\theta f\rangle ,
\label{eq:budget}
\end{equation}
because the advective term is skew-symmetric for $\nabla\!\cdot\!\mathbf{u}=0$
and therefore conserves $\langle\theta^2\rangle$ exactly. Note the consequence,
which matters for scheme design: the \emph{nonlinear} (bilinear) transport term
is the conservative one, and the \emph{linear} diffusive term is what dissipates.
The same structure holds for Navier--Stokes, where $\mathbf{u}\!\cdot\!\nabla\mathbf{u}$
conserves kinetic energy and viscosity dissipates it.

When $\nabla\!\cdot\!\mathbf{u}=0$ the transport term also admits the
\emph{conservative} (flux) form
\begin{equation}
\mathbf{u}\!\cdot\!\nabla\theta \;=\; \nabla\!\cdot\!(\mathbf{u}\,\theta),
\label{eq:cons}
\end{equation}
which is equivalent analytically but not after discretisation: the flux form
respects the variance budget far better and is the usual choice.

\subsection{Cases solved so far}

Work to date has deliberately used constant-coefficient reductions of
Eq.~\eqref{eq:ps}, because they have exactly known solutions and therefore
isolate compression error from every other error source:
\begin{align}
\partial_t\theta &= \kappa\nabla^2\theta && \text{pure diffusion}\\
\partial_t\theta + \mathbf{c}\!\cdot\!\nabla\theta &= 0 && \text{rigid translation, } \mathbf{c}\text{ constant}\\
\partial_t\theta + \mathbf{c}\!\cdot\!\nabla\theta &= \kappa\nabla^2\theta && \text{constant-velocity advection--diffusion}
\end{align}
A genuine, space-dependent velocity has since been added: the steady cellular
flow obtained from the streamfunction $\psi = \sin x \sin y$, i.e.
$\mathbf{u} = (\sin x\cos y,\ -\cos x \sin y)$, run both decaying and weakly
forced. The simpler choice $\mathbf{u} = (\sin y,\ \sin x)$ was tried first and
discarded: it is incompressible but degenerate, in that the advective and
conservative discretisations coincide identically, so it cannot distinguish the
transport forms at all.

The velocity is built from $\psi$ with the solver's own stencils,
$\mathbf{u} = (\mathsf{D}_y\psi,\ -\mathsf{D}_x\psi)$, which makes it
divergence-free to round-off at the \emph{discrete} level since $\mathsf{D}_x$
and $\mathsf{D}_y$ act on different axes and commute exactly. (For separable
trigonometric flows analytic sampling happens to be discretely divergence-free
too; the streamfunction route is the general guarantee, not an improvement in
that special case.)

\subsection{Standard ways to solve these}

\begin{itemize}
\item \textbf{Pseudo-spectral} is the reference method for periodic 2D
      turbulence: derivatives are exact in Fourier space, the nonlinear product
      is formed in physical space with $2/3$ dealiasing, and diffusion is handled
      by an integrating factor or exponential time differencing so it imposes no
      timestep restriction. This is how the DNS database underlying the paper was
      produced.
\item \textbf{Finite difference / finite volume} with upwind or WENO
      reconstruction for the transport term, central stencils for diffusion, and
      SSP Runge--Kutta in time. Cheaper per degree of freedom, more dissipative,
      and much easier to express as a local operator.
\item \textbf{Stiffness.} Explicit diffusion requires
      $\Delta t \le \Delta x^2/4\kappa$ in 2D, which becomes punishing on fine
      grids; implicit or exponential treatment of the diffusive term is standard.
\end{itemize}

\paragraph{What we do instead.} The field is stored as a hybrid TT over the
interleaved (Morton) bit ordering, and every spatial operator is a
matrix-product operator acting on it. Stencils are finite differences, so the
semi-discrete system is
\begin{equation}
\frac{d\boldsymbol{\theta}}{dt} \;=\; \mathsf{A}\,\boldsymbol{\theta},
\qquad
\mathsf{A} \;=\; -c_x\mathsf{D}_x - c_y\mathsf{D}_y + \kappa\,\mathsf{L},
\end{equation}
with $\mathsf{A}$ collapsing to a \emph{single} low-bond MPO in the
constant-coefficient case. Time stepping is explicit Runge--Kutta, and the state
is recompressed after every stage. Compression is the entire source of novelty
and the entire source of risk.

\section{The encoding}

Writing the hybrid ansatz out,
\begin{equation}
\theta(x,y) \;=\; \sum_{\alpha=1}^{\chi} A^{0}[B,\alpha]\;\Phi_\alpha(\text{fine coordinates}),
\end{equation}
with $B$ the coarse-block label carried by the dense head and $\Phi_\alpha$ the
fine-scale shape functions carried by the compressed tail. So the hybrid TT is a
\textbf{block-wise reduced basis with a globally shared dictionary}. Two
consequences:
\begin{itemize}
\item The coarse-grained field is a linear functional of the head alone, so
      ``the large scales live in the head'' is literally true.
\item The large scales carry no rank error. For advection-dominated dynamics this
      is the central argument for the encoding: exact large-scale transport with
      rank-compressed small-scale structure.
\item The same framing predicts the \emph{seam}: neighbouring blocks get
      independent amplitude vectors and nothing enforces continuity across a
      block boundary.
\end{itemize}

\paragraph{Conventions.} $L$ = bits per axis (grid $2^L\times2^L$); $2L$ total
interleaved legs, even legs carrying rows ($y$) and odd legs columns ($x$); $N$ =
legs held by the dense head, so the head spans $2^{N/2}\times2^{N/2}$ coarse
blocks and the seam period is $2^{L-N/2}$ --- \emph{not} $2^{N/2}$, which
coincides only at $N=L$.

\section{What was built}

All of it sits on top of \code{tensnet} via its public API. No \code{tensnet}
source is modified, and --- deliberately stricter --- no \code{tensnet} global
state is read or written, since a wrong global would not raise but would silently
produce wrong physics.

\begin{itemize}
\item \code{hybrid\_ops.py} --- hybrid TT algebra: construction and contraction,
      truncated SVD with a genuine amplitude tolerance, two-sweep compression
      with \emph{per-bond} caps, exact arithmetic (add / scale / axpy /
      Hadamard), MPO application by un-fusing the head, application through an
      explicit fused $2^N\times2^N$ head operator, rank-1 mass projection, and
      pointwise evaluation without contracting the TT.
\item \code{hybrid\_operators.py} --- named stencils, derivative and Laplacian
      MPOs, exact translation, the combined advection--diffusion operator, exact
      MPO algebra, and the isotropic seam smoother.
\item \code{hybrid\_evolve.py} --- Euler / RK2 / SSPRK3 / RK4, a truncation
      policy with optional loss accounting, a stable-timestep helper enforcing
      both the advective and diffusive limits, variable-coefficient transport in
      advective / conservative / skew form with optional forcing, velocity from a
      streamfunction, dense finite-difference and analytic spectral references,
      and diagnostics including a variance-budget measure of numerical
      dissipation.
\item Four test suites, $\sim$117 assertions, all against dense references.
\end{itemize}

All of this lives in \code{time\_evolution/}, deliberately separated from the
paper material in the repository root. The dependency runs one way only: these
modules read \code{config.py} and \code{hybrid\_MPS.py} from the root, and
nothing in the root imports from them, so the paper figures reproduce without
this directory.

The seam smoother is the isotropic repair --- the \code{\_int64\_xy} treatment,
not the $x$-only one. With $P_{x,y}$ the seam projectors, $Q=I-P$, and $A_{x,y}$
the periodic two-neighbour averages,
\begin{equation}
S \;=\; Q_x Q_y \;+\; P_x Q_y A_x \;+\; Q_x P_y A_y \;+\; P_x P_y \tfrac{1}{2}(A_x+A_y).
\end{equation}

\section{What has been checked}

Everything below is verified numerically against a dense reference, not asserted
from reading source.

\begin{itemize}
\item Lossless hybrid round trip; agreement between our constructor and the
      legacy one; the index layout and seam geometry.
\item All arithmetic is exact, with no hidden truncation.
\item Compression: lossless at full rank, monotone in the cap, head physical
      dimension preserved, head bond capped independently of the tail bonds, and
      norm genuinely lost when truncating.
\item Every operator --- first and second derivatives at 2nd and 4th order, upwind
      stencils, Laplacian, translation, seam smoother --- against dense, on both
      axes, at several $(L,N)$.
\item Operator composition, and agreement between composing then applying versus
      applying twice.
\item The seam smoother reproduces the dense reference implementation exactly, is
      isotropic under transposition, and touches only seam lines.
\item All four time steppers reproduce their dense counterparts.
\item The fused-head and un-fused-head paths agree to round-off, both as single
      applications and over multi-step evolutions.
\item Pointwise evaluation agrees with the dense field.
\item The variable-coefficient right-hand side against dense, for all three
      transport forms, with and without diffusion and forcing.
\item Variance conservation of the three transport forms, at several
      resolutions and timestep sizes.
\end{itemize}

\paragraph{Verified \code{tensnet} traps.} None of these raise; all silently
corrupt. The global 2D architecture defaults to the non-interleaved layout;
\code{get\_diff\_mpo} defaults to 8th-order stencils; MPS addition and
multiplication silently compress at a global $l_2$ tolerance; \code{tt.svd}
treats its \code{eps} as a squared-energy fraction, so the achievable amplitude
error saturates around $\sqrt{\epsilon}$; \code{tt.svd} renormalises singular
values after truncation by default; \code{shift\_mpo} has inverted flags and is
non-periodic; \code{MPO.\_\_matmul\_\_} and \code{MPO.contract()} cannot handle
hybrid shapes at all.

\section{Findings and remarks}

\begin{itemize}
\item \textbf{The rank saturates in a genuine cascade.} This is the central
      result. In the cellular flow the intrinsic rank of the \emph{dense}
      solution --- a property of the physics, not of the scheme --- rises during the
      initial stretching and then plateaus for the rest of the run rather than
      growing without bound. An adaptive-tolerance TT tracks it a few units
      below throughout, with the error flat over a thousand steps. Constant-velocity
      transport could not test this; a real cascade can, and the encoding holds.
\item \textbf{Error saturates at fixed rank too}, at every bond dimension
      tried, reproducing what was already seen for pure translation.
\item \textbf{Bulk statistics survive rank starvation far better than
      pointwise accuracy.} At a third of the rank the field actually needs, the
      pointwise error is a few percent while the variance still tracks the
      reference to about one part in a thousand. Encouraging for a
      statistics-focused study, and worth testing directly against the paper's
      diagnostics.
\item \textbf{Numerical dissipation is negligible against molecular
      dissipation} --- a fraction of a percent of it, even at a third of the
      required rank. On this evidence plain truncation is adequate and a
      norm-conserving integrator is not needed, which retires the main open
      question about the scheme.
\item \textbf{The skew-symmetric transport form is essential.} The advective
      and conservative discretisations drift in variance by equal and opposite
      amounts, and their average is conserved to machine zero --- at every
      resolution, timestep and velocity construction tested. The drift of the
      one-sided forms grows sharply as the grid coarsens or the run lengthens,
      so it is not a small effect.
\item \textbf{Centred advection constrains the time scheme.} The spatial
      operator has purely imaginary eigenvalues, so forward Euler is outright
      unstable and RK2 is \emph{latently} unstable: its amplification factor on
      the imaginary axis exceeds one, which stays invisible on a smooth field and
      would bite precisely when the cascade generates high-wavenumber content.
      SSPRK3 or RK4 is required.
\item \textbf{Performance came almost entirely from one change.} Compression
      sweeps dominate the cost, and on the matrix shapes they produce torch's CPU
      LAPACK path is several times slower than scipy's; switching the sweeps gave
      roughly an order of magnitude end to end. Thread count mattered before the
      switch and does not after. Two further optimisations that flop counting
      predicted would help --- replacing an einsum with a matmul, and cutting the
      number of truncations per step --- both measured within noise, because
      fewer-but-larger compressions offset the reduced call count.
\item \textbf{The low-rank structure survives advection.} A one-cell translation
      inflates the state rank by of order ten percent, transiently, and returns it
      \emph{exactly} to baseline whenever the shift is a whole number of coarse
      blocks. Verified at two different $(L,N)$. An adaptive-tolerance scheme
      tracks the intrinsic rank of the exact solution to within one or two, so
      the method accumulates essentially no spurious rank.
\item \textbf{Truncation error saturates rather than compounding.} Under a hard
      rank cap over hundreds of advection steps the error stops growing and the
      field settles onto a representable state that then simply translates. This
      was the principal risk identified at the outset and it did not materialise
      for constant-velocity transport.
\item \textbf{Lossy truncation per step is a property of the cap, not of
      advection.} It appears when the rank cap sits below the field's intrinsic
      rank --- which is exactly the fixed-compression-ratio regime the study works
      in, and a legitimate choice, but the loss should be attributed correctly.
\item \textbf{Rank capping and tolerance capping answer different questions.} A
      cap fixes the compression ratio and lets error float; a tolerance fixes
      error and lets rank float, revealing what the physics actually demands.
      Both are supported and both are needed.
\item \textbf{The seam smoother should not run inside the time loop.} It is a
      low-pass filter and is not idempotent, so iterating it is diffusion. It
      never beat leaving it off. The \emph{severity} is configuration dependent
      --- it scales with the fraction of cells on seam lines, $\approx 2^{1-L+N/2}$
      --- so the earlier dramatic figures came from an unrepresentatively small
      block width. The seam is an artifact of the representation, so its repair
      belongs on the output being analysed, not on the state being evolved.
\item \textbf{Mass conservation is cheap and worth having.} Truncation drifts the
      mean; a rank-1 constant restores it to round-off at no accuracy cost,
      provided it is applied \emph{after} truncation. Unlike rescaling the $l_2$
      norm it does not move variance between scales.
\item \textbf{Do not renormalise after truncation.} Truncation is an orthogonal
      projection and losing norm is a theorem, not a defect. Rescaling puts the
      discarded small-scale variance back into the retained large scales --- a
      spurious inverse cascade --- and, worse, disables the very diagnostic that
      detects under-resolution.
\item \textbf{Consequence for the published data.} Because rescaling does not
      change the retained subspaces, the legacy pipeline's output is exactly
      plain truncation times a \emph{global scalar} that tends to unity as the
      bond dimension grows. Every scale-invariant statistic --- flatness,
      normalised structure-function ratios, normalised increment PDFs --- is
      therefore unchanged; only the absolute level of $E(k)$ can have moved. This
      still needs confirming on the real DNS snapshots.
\item \textbf{The fused head does not scale.} Single-stencil operators have very
      sparse fused heads, but operators that are sums of terms --- the Laplacian,
      the seam smoother --- have essentially dense ones, so sparsity buys nothing
      for exactly the operators one most wants. It is modestly faster at small
      $N$ but becomes infeasible in memory at production $N$; the un-fusing path,
      whose cost carries no $4^N$ term, is the one that scales.
\item \textbf{The cascade run was decaying, not statistically steady.} The
      forcing amplitude was set to balance the \emph{initial} dissipation, but
      once the cascade steepens gradients the dissipation grows by orders of
      magnitude, so the field decays anyway. That the rank plateau holds while
      the variance falls by two orders is arguably a stronger statement, but the
      statistically-steady state is the regime the paper's statistics describe
      and it has not yet been reached.
\item \textbf{Traps for the experimenter.} Rank experiments need a field
      whose rank has room to move: a multiscale field is full rank and the
      measurement merely saturates the chain geometry. A total error that falls
      \emph{below} the pure finite-difference floor is coincidental error
      cancellation, not an improvement. A diagnostic that aggregates truncation
      losses across a whole right-hand side does not measure dissipation of the
      state. And a transport test flow must be genuinely two-dimensional, or the
      discretisation forms it is meant to separate coincide identically.
\end{itemize}

\section{Corrections to earlier statements}

Recorded because the reasoning matters more than the conclusions:
\begin{itemize}
\item ``Translation is rank-preserving'' (wrong), then ``translation is not
      rank-preserving, so every step must be lossy'' (also wrong, over-corrected).
      The measured behaviour is the first bullet of the previous section.
\item ``Apply the seam smoother once per step, inside the loop.'' It is harmful
      at every cadence tested.
\item ``The nonlinear term is the norm-changing one.'' For a passive scalar with
      $\nabla\!\cdot\!\mathbf{u}=0$ the transport term conserves variance exactly
      and the linear diffusive term dissipates; see Eq.~\eqref{eq:budget}.
\item \textbf{The truncation-loss diagnostic was wrong by four orders of
      magnitude}, and in the alarming direction: it reported numerical
      dissipation at roughly seventeen times the molecular value when the true
      figure is about two thousandths of it. It summed every compression,
      including the dozen or so per right-hand-side evaluation performed on
      intermediates whose norms are those of the right-hand side rather than of
      $\theta$. Caught only because it contradicted the variance history, which
      tracked the reference closely. Replaced by a variance-budget measure, with
      a warning left on the raw accumulator.
\item \textbf{Two performance predictions from flop counting, both wrong.}
      Replacing an einsum with a matmul and cutting the truncation count per step
      were each expected to help and each measured within noise. The actual win
      was in a place flop counting does not look at all: the LAPACK backend.
      Relatedly, the compression granularity inside the transport term was
      argued in both directions before being settled by measurement.
\item \textbf{Two wrong diagnoses of a changed variance drift.} A discrepancy
      against an earlier measurement was attributed first to the timestep and
      then to the velocity construction; both were ruled out by direct test, and
      the cause was the grid resolution together with the run length implied by a
      fixed step count.
\end{itemize}

\section{Status and next steps}

\paragraph{Done.} The validated baseline is committed and separated from the
paper material. The cellular-flow cascade has been run and answers the central
question: the rank saturates. The dissipation comparison, which came with it,
shows truncation dissipation to be negligible against molecular diffusion, so the
conditional TDVP upgrade is not currently warranted. Compression performance has
been profiled and improved by about an order of magnitude.

\begin{enumerate}[leftmargin=*]
\item \textbf{A properly forced steady state.} The cascade run so far decays;
      the forcing needs to be far stronger to reach the statistically steady
      regime the paper's statistics describe. Cheap now that a run costs minutes
      rather than hours. \emph{Settles:} whether the rank plateau survives in a
      sustained cascade rather than a decaying one.

\item \textbf{Hybrid versus standard TT at matched compression.} The actual
      scientific claim and the true sequel to the paper: does the hybrid encoding
      reproduce intermittency statistics \emph{under evolution} better than a
      standard TT at equal $\rho$? Needs a standard-TT evolution path, the
      existing degree-of-freedom matching, and the paper's diagnostics recast as
      functions of time. Compression quality does not imply dynamical quality;
      this is what measures it. The observation that variance survives rank
      starvation far better than pointwise accuracy makes the outcome genuinely
      uncertain, and therefore worth measuring.

\item \textbf{Scale to the production geometry} ($4096^2$). Use the un-fusing
      head path: the fused head is dense for summed operators and does not scale.
      Two further compression optimisations remain in reserve --- fusing operator
      application with truncation, and tracking the canonical gauge to skip a
      sweep --- worth perhaps another factor of two, and more relevant at this
      size than at $128^2$.

\item \textbf{Real DNS velocity fields}, together with the deferred check of the
      renormalisation effect on the published figures. Both need the dataset.

\item \textbf{Conditional, and currently not indicated: TDVP or implicit
      stepping.} The dissipation measurement argues against the former for now;
      the latter becomes necessary only if the diffusive timestep restriction
      starts to bite at production resolution.
\end{enumerate}

\paragraph{Beyond the passive scalar.} The natural extension is 2D
Navier--Stokes in vorticity form,
$\partial_t\omega + \mathbf{u}\!\cdot\!\nabla\omega = \nu\nabla^2\omega$ with
$\mathbf{u}$ recovered from $\omega$ by an elliptic solve, which adds a Poisson
inversion --- a natural fit for the fused/sparse head --- and a second quadratic
invariant, since 2D conserves both energy and enstrophy and truncation will
damage them at different rates.

\end{document}
